% from ega4-2-fr.tex
\oldpage[IV-2]{161}

En vertu de ce lemme, on peut supposer qu'il existe dans $X$ un
sous-préschéma intègre $Y$ de point générique $y$ tel que $\mathcal F$
vérifie $(S_n)$ au point $y$ et en tous les points de $X-Y$, et tout se
borne à vérifier qu'il existe dans $X$ un voisinage ouvert de $y$ tel que
$\mathcal F$ vérifie $(S_n)$ dans ce voisinage. Distinguons alors deux cas:

$1^\circ$ $y$ est point maximal de $X$; comme il existe un voisinage ouvert
de $y$ ne rencontrant aucune composante irréductible de $X$ autre que
$\overline{\{y\}}$, on peut supposer que $X$ est irréductible, donc a même
espace sous-jacent que $Y$, de sorte que $Y$ est défini par le Nilradical de
$\mathcal O_X$, qui est nilpotent. D'autre part, on peut, en remplaçant $X$
par un voisinage ouvert de $y$, supposer que $\mathcal F$ est normalement
plat le long de $Y$ (6.9.1); il résulte alors de (6.10.4) que
$U_{S_n}(\mathcal F)=U_{S_n}(\mathcal O_Y)$, et comme ce dernier est par
hypothèse un voisinage de $y$ dans $X$, cela prouve la proposition dans ce
cas.

$2^\circ$ $y$ n'est pas un point maximal de $X$, autrement dit (puisque
$\operatorname{Supp}(\mathcal F)=X$), $\dim(\mathcal F_y)\geq1$, donc aussi
$\operatorname{prof}(\mathcal F_y)\geq1$ puisque $\mathcal F$ vérifie par
hypothèse $(S_n)$ (et a fortiori $(S_1)$) au point $y$. Remplaçant au besoin
$X$ par un voisinage ouvert de $y$, on peut donc supposer qu'il existe une
section $f$ de $\mathcal O_X$ au-dessus de $X$, $\mathcal F$-régulière et
telle que $f_y\in\mathfrak m_y$, ou encore $f(y)=0$ (0, 15.2.4); on a donc
encore $f(x)=0$ pour tout $x\in Y$ (0$_{\rm I}$, 5.5.2). On sait que
$\mathcal F/f\mathcal F$ vérifie $(S_{n-1})$ au point $y$ (5.7.6).
Appliquant l'hypothèse de récurrence, et remplaçant au besoin $X$ par un
voisinage ouvert de $y$, on peut donc supposer que
$\mathcal F/f\mathcal F$ vérifie $(S_{n-1})$ en tout point de $X$. Mais pour
tout $x\in Y$, la relation $f(x)=0$ entraîne que l'on a
$\operatorname{prof}(\mathcal F_x/f_x\mathcal F_x)=
\operatorname{prof}(\mathcal F_x)-1$ et
$\dim(\mathcal F_x/f_x\mathcal F_x)=\dim(\mathcal F_x)-1$
(0, 16.3.4 et 16.4.6); la relation
\[
\operatorname{prof}(\mathcal F_x/f_x\mathcal F_x)\geq
\inf(n-1,\dim(\mathcal F_x/f_x\mathcal F_x))
\]
est donc équivalente à
\[
\operatorname{prof}(\mathcal F_x)\geq\inf(n,\dim(\mathcal F_x)).
\]
Comme on a supposé que $\mathcal F$ vérifie $(S_n)$ en tout point de $X-Y$,
cela termine la démonstration.
\end{proof}

\begin{corollary}[6.11.7]
\label{IV.2.s6.11.7-fr}
Les notations étant celles de (6.11.6):
\begin{enumerate}
\item[(i)] L'ensemble $U_{S_1}(\mathcal F)$ est ouvert dans $X$.
\item[(ii)] Pour que l'ensemble $U_{S_2}(\mathcal F)$ soit ouvert il suffit
que tout point maximal $x$ de $\operatorname{Supp}(\mathcal F)$,
appartenant à $U_{S_2}(\mathcal F)$, soit intérieur à
$U_{S_2}(\mathcal F)$.
\end{enumerate}
\end{corollary}

\begin{proof}
L'assertion (i) a été prouvée dans le cours de la démonstration de (6.11.6);
d'autre part, pour $n=2$ le cas $2^\circ$ de la démonstration de (6.11.6)
est valable sans aucune hypothèse sur $X$, puisque (avec les mêmes
notations) $U_{S_1}(\mathcal F)$ et
$U_{S_1}(\mathcal F/f\mathcal F)$ sont ouverts dans $X$. Quant au cas
$1^\circ$ de cette démonstration, l'hypothèse assure précisément qu'il est
inutile de le considérer.
\end{proof}

